\documentclass[10pt,a4paper]{article}
\usepackage{amsmath,amssymb,amsthm,graphicx}
\usepackage{amsmath,amssymb,amsthm,hyperref,xfrac}
\providecommand{\ie}{i.e.\ }
\providecommand{\skpt}{\smallskip}
\usepackage[scr=rsfs,cal=euler]{mathalfa}
\multlinegap0pt

\hyphenation{begin decom-po-si-tion sequence}
\newenvironment{enumeratea}
{\bgroup\def\theenumi{\alph{enumi}}\begin{enumerate}}
{\end{enumerate}\egroup}

\newcommand{\ba}{{\boldsymbol a}}
\newcommand{\bb}{{\boldsymbol b}}
\newcommand{\bd}{{\boldsymbol d}}
\newcommand{\bu}{{\boldsymbol u}}
\newcommand{\bx}{{\boldsymbol x}}
\newcommand{\bmy}{{\boldsymbol y}}

\newcommand{\bD}{{\boldsymbol D}}
\newcommand{\bL}{{\boldsymbol L}}
\newcommand{\bP}{{\boldsymbol P}}
\newcommand{\bV}{{\boldsymbol V}}
\newcommand{\bX}{{\boldsymbol X}}
\newcommand{\bZ}{{\boldsymbol Z}}

\newcommand{\Psfrac}[2]{(\sfrac{#1}{#2})}
\newcommand{\psfrac}[2]{\sfrac{(#1)}{#2}}
\newcommand{\bigO}[1]{O(#1)}

\usepackage{mathtools}
\usepackage{bbm}

\newcommand{\NN}{\mathbb{N}}
\newcommand{\ZZ}{\mathbb{Z}}
\newcommand{\RR}{\mathbb{R}}
\newcommand{\CC}{\mathbb{C}}
\newcommand{\PP}{\mathbb{P}}
\newcommand{\EE}{\mathbb{E}}

\newcommand{\boldalpha}{\boldsymbol{\alpha}}
\newcommand{\bolddelta}{\boldsymbol{\delta}}
\newcommand{\boldrho}{\boldsymbol{\rho}}
\newcommand{\dee}{\,\mathrm{d}}
\newcommand{\one}{\mathbbm{1}}
\newcommand{\BC}{\mathcal{B}}
\newcommand{\EC}{\mathcal{E}}
\newcommand{\RS}{\mathscr{R}}
\newcommand{\prob}[1]{\PP[ #1 ]}
\newcommand{\bigoo}[1]{O\!\left(#1 \right)}
\newcommand{\abs}[1]{\left\lvert #1 \right\rvert}
\newcommand{\floor}[1]{\left\lfloor #1 \right\rfloor}
\newcommand{\riemint}[4]{\int_{#1}^{#2}\! #3 \, \mathrm{d}#4}
\newcommand{\maxnorm}[1]{\left\| #1 \right\|_{\infty}}

\renewcommand{\Re}{\textup{Re}}
\renewcommand{\Im}{\textup{Im}}

\newcommand{\dtv}{\mathrm{d}_{\mathrm{TV}}}
\newcommand{\dw}{\mathrm{d}_{\mathrm W}}

\let\epsilon\varepsilon

\newcommand{\pextra}{P_{\mathtt{extra}}}
\newcommand{\parratia}{P_{\mathtt{Arratia}}}

\newcommand{\simplex}{\Delta^{k-1}}

\newcommand{\bg}{\big}
\newcommand{\bgg}{\Big}
\newcommand{\bggg}{\bigg}

\newcommand{\Rnt}{R_{\mathrm{NT}}}
\newcommand{\Rprob}{R_{\mathrm{PD}}}
\newcommand{\ECnt}{\EC_{\mathrm{NT}}}
\newcommand{\ECprob}{\EC_{\mathrm{PD}}}

\newcommand{\Bnt}{\mathcal{B}_{\mathrm{NT}}}
\newcommand{\Bprob}{\mathcal{B}_{\mathrm{PD}}}

\DeclareMathOperator{\vard}{Var}
\DeclareMathOperator{\Dir}{Dir}

\theoremstyle{plain}
\newtheorem{cor}{Corollary}[section]
\newtheorem{lem}[cor]{Lemma}
\newtheorem{prop}[cor]{Proposition}
\newtheorem{conj}[cor]{Conjecture}
\newtheorem{theo}{Theorem}

\theoremstyle{definition}
\newtheorem{example}{Example}
\newtheorem{deff}[cor]{Definition}
\newtheorem{rem}{Remark}[section]
\newtheorem*{rems}{Remark}
\newtheorem*{remss}{Remarks}

\makeatletter
\ams@newcommand{\multiint}[1]{\DOTSI\protect\MultiIntegral{#1}}
\renewcommand{\MultiIntegral}[1]{%
\edef\ints@c{\noexpand\intop
\ifnum#1=\z@\noexpand\intdots@\else\noexpand\intkern@\fi
\replicate{#1-2}{\noexpand\intop\noexpand\intkern@}%
\noexpand\intop
\noexpand\ilimits@
}%
\futurelet\@let@token\ints@a
}
\makeatother
\ExplSyntaxOn
\cs_new:Npn \replicate #1 #2 { \prg_replicate:nn { #1 } { #2 } }
\ExplSyntaxOff
\usepackage{cleveref}
\crefname{originalTheorem}{original Theorem}{original Theorems}
\Crefname{originalTheorem}{Original Theorem}{Original Theorems}
\makeatletter
\newlabel{thm:fact into k parts}{{2}{source}{}{sourcepointer-1}{}}
\newlabel{thm:fact into k parts@cref}{{[originalTheorem][0][]2}{[1][0][]source}{}{}{}}
\makeatother
\title{Uniform boundary probability for multiplicatively weighted random integer factorizations}
\author{Tony Haddad, Dimitris Koukoulopoulos}
\date{doi:10.5802/jep.317; CC BY4.0}
\begin{document}
\maketitle
\noindent\textit{Editorial note.} The counted item is the number-theoretic boundary-event lemma, not the separate PD-radius lemma or the old coupling theorem.\par\medskip
\noindent\textit{Editorial note.} The constants sought are uniform in alpha\_i,u,k and the permitted allocation law f; the two complementary hyperbola regions are essential for retaining that scope.\par\medskip
\noindent\textit{Editorial note.} The first mean-value estimate1682 is a named standard prerequisite. Its local weighted harmonic-window consequence1683 and proof1687-1700 remain in the solution chain.\par\medskip
\noindent\textit{Editorial note.} At source1699 the numerical intermediate count of the logarithmic partition intervals is understated: the defining interval1696 has logarithmic length1+2log y. The proof uses only O(log y), which is also stated on1699 and is valid uniformly fory>=2. The original author line remains unchanged; the implicit absolute constant absorbs this numerical count.\par\medskip
\noindent\textit{Editorial scope.} One integer-factorization boundary bound. Its square-full moment and weighted harmonic-window auxiliaries are not separate counted items. The full Dirichlet-law conclusion is not counted on top of these prepared estimates.\par\medskip
\noindent\textit{Allowed standard prerequisites (editorial).}\begin{itemize}
\item Mean-value upper bound for alpha to the number of distinct prime factors. Uniformly for alpha in[0,1] and y>=2, sum\_(m<=y) alpha\textasciicircum{}(omega(m)) <= C y(log y)\textasciicircum{}(alpha-1).. Bound one factor after the proof-specific square-free decomposition and hyperbola split. The short harmonic window bound is proved locally.
\item Classical Mertens theorem and partial summation. Prove the local weighted harmonic-window estimate1678-1700.
\item Counting square-full integers. There are O(sqrt z) square-full integers at most z; each can be written a\textasciicircum{}2b\textasciicircum{}3 with b square-free, and sum\_b b\textasciicircum{}(-3/2) converges.. Bound the square-full allocation count and the moment sum. No boundary probability estimate is supplied.
\end{itemize}\medskip
\ifdefined\hypertarget\hypertarget{sourcepointer-1}{}\fi
\noindent\textit{Source reference.} Original Theorem 2: sources/arratia-dirichlet-jep-2025-317/source0.tex, source lines 349--359. The definition footnote says the allocation conditions will be used for this later theorem; that theorem is not used to prove the present integer boundary estimate.\par\medskip
Let $N_x$ be an integer chosen uniformly at random from the set $\ZZ \cap [1, x]$. We~may then factor it uniquely as $N_x = P_1P_2 \cdots$ with the $P_i$'s forming a non-increasing sequence of primes or ones. In~1972, Billingsley \cite{Bill72} showed that, for any fixed positive integer $r$, the joint distribution of the random vector

Let $\simplex$ be the set of $k$-tuples $\boldalpha=(\alpha_1, \ldots, \alpha_k) \in \RR^k_{\ge 0}$ satisfying $\alpha_1+\cdots+\alpha_k=1$. We~also need to define a special class of functions:

\begin{deff}[The class of functions $\mathcal F_k(\boldalpha)$]\label{dfn:Fktheta}
Given $k\in\ZZ_{\ge2}$ and $\boldalpha\in \simplex$, let $\mathcal F_k(\boldalpha)$ be the set of functions $f \colon \NN^{k} \to \RR_{ \ge 0}$ satisfying the following three properties:
\begin{enumeratea}
\item\label{dfn:Fkthetaa} We have $\sum_{d_1\cdots d_k = n} f(d_1, \ldots, d_k) = 1$ for all $n \in \NN$. Thus, for each fixed $n \in \NN$, the function $f(\bd)$ defines a probability mass function on the set of vectors $\bd \in \NN^k$ satisfying $d_1\cdots d_k = n$.

\item\label{dfn:Fkthetab} Whenever $\bd$ satisfies
\[
d_i \coloneqq \begin{cases}
p &\text{if $i = j$}, \\
1 &\text{if $i \ne j$}.
\end{cases}
\]
for some $1 \le j \le k$ and prime $p$, then $f(\bd) = \alpha_j$.

\item\label{dfn:Fkthetac} The function $f$ is \emph{multiplicative}, \ie for any vectors $\ba, \bb \in \NN^k$ with $a_1\cdots a_k$ and $b_1 \cdots b_k$ being co-prime, we~have the property
\[
f(a_1b_1, \ldots, a_kb_k) = f(a_1, \ldots, a_k) \cdot f(b_1, \ldots, b_k).
\]
\end{enumeratea}
\end{deff}

\begin{rems}
Let $\omega(n)$ denote the number of distinct prime factors of $n$. If~$d_1\cdots d_k$ is a square-free number and $f\in\mathcal{F}_k (\boldalpha)$, then properties \eqref{dfn:Fkthetab} and \eqref{dfn:Fkthetac} of the definition above imply that
\begin{equation}
\label{eq:f on squarefrees}
f(d_1,\dots,d_k) = \prod_{j=1}^k \alpha_j^{\omega(d_j)}.
\end{equation}
\end{rems}

We will use the class of functions $\mathcal{F}_k(\boldalpha)$ to define certain ``random factorizations'' into $k$ parts of a random integer. Specifically, let us fix $f\in\mathcal{F}_k(\boldalpha)$ and $x\ge1$. We~then define a {\it random $k$-factorization corresponding to $f$} to be a random vector $\bD_{f, x}=(D_{f,x,1},\dots,D_{f,x,k})$ taking values on $\NN^k$ and satisfying the formula\footnote{Strictly speaking, we~only need property \eqref{dfn:Fkthetaa} of Definition \ref{dfn:Fktheta} to define $\bD_{f,x}$. But we will also need the other two properties when proving Theorem \ref{thm:fact into k parts} below.}
\begin{equation}
\label{eq:prob mass function f}
\PP\bgg[D_{f,x,i} = d_i \ \, \forall i \le k\, \bgg|\, N_x = n\bgg] = f(d_1, \ldots, d_k)
\end{equation}
for all $n\in\NN$ and all $k$-tuples $(d_1,\dots,d_k)\in\NN^k$ with $d_1\cdots d_k=n$.

\subsection{Notation}
We let $\log_j$ denote the $j$-iteration of the natural logarithm, meaning that $\log_1=\log$ and $\log_j=\log\circ \log_{j-1}$ for $j\ge2$.

Throughout the paper, the letter $p$ is reserved for prime numbers and the letter $n$ is reserved for natural numbers, unless stated otherwise. Given such $p$ and $n$, we~write $\nu_p(n)$ for the $p$-adic valuation of $n$, that is to say the largest integer $v\ge0$ such that $p^v|n$. In~addition, we~write $\omega(n)$ for the number of distinct prime factors of $n$.

Moreover, for $n\in\NN$, we~let $s(n)$ denote its largest square-full divisor. Also, we~let $n^\flat \coloneqq n/s(n)$ and we note that $n^\flat$ is always square-free and co-prime to $s(n)$.

We write $\pi(x)$ for the number of primes $\le x$. We~shall also use heavily Chebyshev's function $\theta(x) \coloneqq \sum_{p\le x}\log p$.

To describe various estimates, we~use Vinogradov's notation $f(x) \ll g(x)$ or Landau's notation $f(x) = \bigoo{g(x)}$ to mean that $\abs{f(x)} \le C\cdot g(x)$ for a positive constant~$C$. If~$C$ depends on a parameter $\alpha$, we~write $f(x) \ll_\alpha g(x)$ or $f(x) = O_\alpha(g(x))$. If~two positive functions $f, g$ have the same order of magnitude in the sense that $f(x) \ll g(x) \ll f(x)$, then we write $f(x) \asymp g(x)$.

If $P$ is some proposition, then the indicator function $\one _P$ will be equal to 1 if $P$ is true and $0$ if $P$ is false.

\[
\ECnt(\simplex_{\mathbf{u}})^c\subseteq\! \bigcup_{\substack{i<k \\ u_i\neq 1}} \Bnt^{(i)}(u_i),
\quad\text{where}\quad
\Bnt^{(i)}(u) \coloneqq \bgg\{ \bg| \log(D_{f,x,i}/x^u)\bg|\le 3\log(x/N_x^\flat) \bgg\},
\]

\begin{prop}[The prime number theorem]
\label{pnt}
For $x \ge 2$, we~have
\[
\theta(x) = x + \bigO{\sfrac{x}{(\log x)^3}}
\]
and
\[
\psi(x) = x+\bigO{\sfrac{x}{(\log x)^3}}.
\]
\end{prop}

We shall also need a stronger version of Mertens's estimate, which in its classical form, yields an estimate for the product $\prod_{p \le x} (1-1/p)$.

\begin{prop}[Strong Mertens's estimate]
\label{prop:mertens}
For $x \ge 2$, we~have
\[
\sum_{p^k \le x} \frac{1}{kp^k} = \log_2 x + \gamma + \bigO{\sfrac{1}{(\log x)^3}}.
\]
\end{prop}

\begin{proof}
Using the prime number theorem (Proposition \ref{pnt}), there exists a constant $c_0$ satisfying
\[
\sum_{p^k \le x} \frac{1}{kp^k} = \int_{2^-}^x \frac{\dee \psi(t)}{t\log t} = \log_2 x + c_0 + \bigO{\sfrac{1}{(\log x)^3}}.
\]
for $x \ge 2$. For the proof of $\gamma = c_0$, see the proof of \cite[Th.\,2.7]{MV07}.
\end{proof}

\begin{lem}
\label{lem:bound of prob of ECnt}
For $i\in\{1,2,\dots,k-1\}$ and $u\in(0,1)$, we~have the uniform estimate
\[
\PP\bg[\Bnt^{(i)}(u) \bg] \ll \frac{1}{(1+ u\log x)^{1-\alpha_i}(1+(1-u)\log x)^{\alpha_i}}.
\]
\end{lem}

\begin{lem}[Properties of $N_x$]
\label{lem:properties of N}
Fix $\alpha\in[0,1)$ and $\beta\in[0,1/2)$. Uniformly over $x\ge1$, we~have
\[
\EE\bg[(x/N_x)^\alpha s(N_x)^\beta\bg] \ll_{\alpha,\beta} 1.
\]
\end{lem}

\begin{proof} We must show that
\[
S\coloneqq \sum_{n \le x} (x/n)^\alpha s(n)^\beta\ll_{\alpha,\beta} x.
\]
Indeed, if we let $b=s(n)$ and $a=n^\flat=n/b$, then
\[
S\le \sum_{\substack{b\le x \\ b\ \text{square-full} }} x^\alpha b^{\beta-\alpha} \sum_{a\le x/b} a^{-\alpha}
\ll_\alpha \sum_{\substack{b\le x \\ b\ \text{square-full} }} x^\alpha b^{\beta-\alpha} (x/b)^{1-\alpha}
\le x \sum_{b\ \text{square-full} } b^{\beta-1},
\]
where we used our assumption that $\alpha<1$.
Since we have assumed that $\beta<1/2$, the last sum over $b$ converges, thus completing the proof.
\end{proof}

\section{Proof of Lemma \ref{lem:bound of prob of ECnt} using elementary number-theoretic techniques}\label{sec:boundary event NT}

Note that
\begin{equation}
\label{eq:small square-free part}
\PP\bg[N_x^\flat\le x/(\log x)^3\bg] \le \frac{1}{\log x} \EE\bg[(x/N_x^\flat)^{1/3}\bg] \ll \frac{1}{\log x}
\end{equation}
by Lemma \ref{lem:properties of N} applied with $\alpha=\beta=1/3$.

Fix now $i\in\{1,\dots,k-1\}$ and $u\in(0,1)$. For simplicity, let us write
\[
\alpha\coloneqq \alpha_i
\]
for the remainder of this section. Given $z$, let us define
\[
S(z)\coloneqq \sum_{\substack{n\le x \\ n^\flat \in(\sfrac{x}{z},\sfrac{2x}{z}]}}\; \sum_{\substack{d_1\cdots d_k = n \\ d_i\in[x^uz^{-3}, x^uz^3] }}\hspace*{-5mm} f(d_1,\dots,d_k).
\]
Using \eqref{eq:small square-free part} and \eqref{eq:prob mass function f}, we~find that
\begin{equation}
\label{eq:bounding Bnt 1}
\PP\bg[\Bnt^{(i)}(u)\bg] \le \frac{1}{\floor{x}} \sum_{1\le m\le 5\log_2x}\hspace*{-3mm} S(2^m) + O\bg(\sfrac{1}{\log x}\bg).
\end{equation}
Hence Lemma \ref{lem:bound of prob of ECnt} follows readily from the estimate below:

\begin{lem}
There exists a universal constant $x_0\ge2$ such that if $x\ge x_0$ and $z\in[2,(\log x)^4]$, then we have the uniform estimate
\[
S(z) \ll \frac{\log z}{\sqrt{z}} \cdot \frac{x}{(1+ u\log x)^{1-\alpha}(1+(1-u)\log x)^{\alpha}}.
\]
\end{lem}

\begin{proof}
If $d_1\cdots d_k=n$, then we may uniquely write $d_j=d_j'd_j''$ with $d_j'|n^\flat$ and $d_j''|s(n)$. In~particular, $d_j''\le s(n)\le x/n^\flat\le z$ for all $j$. Hence, if $d_i\in[x^u z^{-3},x^uz^3]$, we~must also have that $d_i'\in[x^uz^{-4},x^uz^3]$. Hence, using property \eqref{dfn:Fkthetac} of Definition~\ref{dfn:Fktheta}, we~deduce that
\[
\sum_{\substack{d_1\cdots d_k = n \\ d_i\in[x^uz^{-3}, x^uz^3] }} \hspace*{-5mm}f(d_1,\dots,d_k)
\le \sum_{\substack{d_1'\cdots d_k' = n^\flat \\ d_i'\in[x^uz^{-4}, x^uz^3] }}\hspace*{-5mm} f(d_1',\dots,d_k')
\sum_{d_1''\cdots d_k'' = s(n)}\hspace*{-3mm} f(d_1'',\dots,d_k'').
\]
Relation \eqref{eq:f on squarefrees} implies that $f(d_1',\dots,d_k') = \prod_{j=1}^k \alpha_j^{\omega(d_j')}$. Moreover, using property \eqref{dfn:Fkthetaa} of Definition \ref{dfn:Fktheta}, we~find that the sum over the $d_j''$'s equals $1$. In~conclusion,
\begin{align*}
\sum_{\substack{d_1\cdots d_k = n \\ d_i\in[x^uz^{-3}, x^uz^3] }}\hspace*{-4mm} f(d_1,\dots,d_k)
&\le \sum_{\substack{d_1'\cdots d_k' = n^\flat \\ d_i'\in[x^uz^{-4}, x^uz^3] }} \prod_{j=1}^k \alpha_j^{\omega(d_j')} \\
&= \sum_{\substack{d| n^\flat \\ d\in[x^uz^{-4}, x^uz^3] }} \hspace*{-5mm} \alpha^{\omega(d)} (1-\alpha)^{\omega(n^\flat/d)}.
\end{align*}
Let $n^\flat=dm$ and $s(n)=b$. If~$n^\flat\in[\sfrac{x}{z},\sfrac{2x}{z}]$ and $d\in[x^uz^{-4},x^uz^3]$, then $m\in [x^{1-u}z^{-4},2x^{1-u}z^3]$. Hence, we~find that
\begin{equation}
\label{eq:bound in terms of S1 and S2}
\begin{aligned}
S(z) &\le \mathop{\sum\sum}_{\substack{dm \in[\sfrac{x}{z},\sfrac{2x}{z}] \\ d\in[x^uz^{-4}, x^uz^3] \\ m\in[x^{1-u}z^{-4},2x^{1-u}z^3] }}\hspace*{-5mm}
\alpha^{\omega(d)} (1-\alpha)^{\omega(m)} \sum_{\substack{b\le z \\ b\ \text{square-full}}}\hspace*{-5mm} 1\\
&\ll \sqrt{z}\hspace*{-5mm} \mathop{\sum\sum}_{\substack{dm \in[\sfrac{x}{z},\sfrac{2x}{z}] \\ d\in[x^uz^{-4}, x^uz^3] \\ m\in[x^{1-u}z^{-4},2x^{1-u}z^3] }}\hspace*{-4mm} \alpha^{\omega(d)} (1-\alpha)^{\omega(m)}
= \sqrt{z}\cdot (S_1+S_2),
\end{aligned}
\end{equation}
where $S_1$ denotes the double sum over $d$ and $m$ with the additional constraint \hbox{$d\le \sqrt{2x/z}$}, and $S_2$ is the corresponding sum over pairs $(d,m)$ with $d>\sqrt{2x/z}$ and $m\le \sqrt{2x/z}$.

First, we~bound $S_1$. Note that for the conditions on $d$ to be compatible we must have that $u\le 2/3$, provided that $x$ is large enough. Assuming this is the case, we~apply twice Proposition \ref{prel:partial summation bounds} to find that
\begin{equation}\label{eq:bound for S1}
\begin{aligned}
S_1 &\le \sum_{\substack{d\le \sqrt{2x/z} \\ d\in[x^uz^{-4}, x^uz^3] }}\hspace*{-5mm} \alpha^{\omega(d)} \sum_{m\le \sfrac{2x}{dz}} (1-\alpha)^{\omega(m)} \\
&\ll \sum_{\substack{d\le \sqrt{2x/z} \\ d\in[x^uz^{-4}, x^uz^3] }}\hspace*{-5mm} \alpha^{\omega(d)} \cdot \frac{x(\log x)^{-\alpha}}{dz} \\
&\ll \frac{x(\log z)(1+ u\log x)^{\alpha-1}(\log x)^{-\alpha}}{z}.
\end{aligned}
\end{equation}

We bound $S_2$ in a similar manner. For this sum to be non-empty, we~have $u\ge1/3$. If~this is the case, then we have
\begin{equation}
\label{eq:bound for S2}
\begin{aligned}
S_2
&\le \sum_{\substack{m\le \sqrt{2x/z} \\ m\in[x^{1-u}z^{-4}, 2x^{1-u}z^3] }}\hspace*{-5mm} (1-\alpha)^{\omega(m)} \sum_{d\le \sfrac{2x}{mz}} \alpha^{\omega(d)} \\
&\ll \sum_{\substack{m\le \sqrt{2x/z} \\ m\in[x^{1-u}z^{-4}, 2x^{1-u}z^3] }}\hspace*{-5mm} (1-\alpha)^{\omega(m)} \cdot \frac{x(\log x)^{\alpha-1}}{mz} \\
&\ll \frac{x(\log z)(1+ (1-u)\log x)^{-\alpha}(\log x)^{\alpha-1}}{z}.
\end{aligned}
\end{equation}
Combining \eqref{eq:bound in terms of S1 and S2}, \eqref{eq:bound for S1} and \eqref{eq:bound for S2}, while keeping in mind that $S_1=0$ if $u>2/3$ and that $S_2=0$ if $u<1/3$, completes the proof of the lemma,
\end{proof}

\begin{prop}
\label{prel:partial summation bounds}
Uniformly for $\alpha\in[0,1]$ and $x,y\ge2$, we~have
\[
\sum_{n \le x} \alpha^{\omega(n)} \ll x(\log x)^{\alpha-1}
\quad \text{and} \quad \sum_{x/y<n \le xy} \frac{ \alpha^{\omega(n)}}{n} \ll (\log y)(\log(xy))^{\alpha-1}.
\]
\end{prop}

\begin{proof}The first bound follows readily with \cite[Th.\,14.2]{Koukoulo19}. Let us now prove the second one.

When $y\ge \sqrt{x}$, we~have
\[
\sum_{x/y<n \le xy} \frac{ \alpha^{\omega(n)}}{n}\le \sum_{p|n\ \Rightarrow p \le y^3} \frac{ \alpha^{\omega(n)}}{n} = \prod_{p\le y^3} \bgg(1+\frac{\alpha}{p-1}\bgg) \ll (\log y)^\alpha
\]
by the inequality $1+t\le e^t$ and Mertens's theorem (Proposition \ref{prop:mertens}). This completes the proof in this case. Finally, assume that $y\le \sqrt{x}$. We~then have
\[
\sum_{x/y<n \le xy} \frac{ \alpha^{\omega(n)}}{n}
\le \sum_{\substack{m\in\ZZ \\ e^m\in[x/y,exy]}} \sum_{n\in(e^{m-1},e^m]} \frac{ \alpha^{\omega(n)}}{n}
\le \sum_{\substack{m\in\ZZ \\ e^m\in[x/y,exy]}} e^{1-m} \sum_{n\le e^m} \alpha^{\omega(n)}.
\]
For every $m$ as above, the innermost sum is $\ll e^m m^{\alpha-1} \asymp e^m (\log(xy))^{\alpha-1}$ by the first part of the lemma and our assumption that $y\le \sqrt{x}$. Since there are $\le1+\log y\ll \log y$ choices for $m$, the needed estimate follows in this last case too.
\end{proof}

\begin{thebibliography}{99}
\bibitem{Arratia98} R. Arratia - “On the central role of scale invariant Poisson processes on (0,∞)”, in Microsurveys in discrete probability (Princeton, NJ, 1997), DIMACS Ser. Discrete Math. Theoret. Comput. Sci., vol. 41, American Mathematical Society, Providence, RI, 1998, p. 21-41
\bibitem{Arratia02} R. Arratia - “On the amount of dependence in the prime factorization of a uniform random integer”, in Contemporary combinatorics, Bolyai Soc. Math. Stud., vol. 10, János Bolyai Math. Soc., Budapest, 2002, p. 29-91
\bibitem{ABT06} R. Arratia, A. D. Barbour \& S. Tavaré - “A tale of three couplings: Poisson-Dirichlet and GEM approximations for random permutations”, Combin. Probab. Comput. 15 (2006) no. 1-2, p. 31-62
\bibitem{BareikisManstavicius07} G. Bareikis \& E. Manstavičius - “On the DDT theorem”, Acta Arith. 126 (2007) no. 2, p. 155-168
\bibitem{Bill72} P. Billingsley - “On the distribution of large prime divisors”, Period. Math. Hungar. 2 (1972), p. 283-289
\bibitem{delaBretecheTenenbaum16} R. de la Bretèche \& G. Tenenbaum - “Sur les processus arithmétiques liés aux diviseurs”, Adv. in Appl. Probab. 48 (2016) no. A, p. 63-76
\bibitem{DDT79} J.-M. Deshouillers, F. Dress \& G. Tenenbaum - “Lois de répartition des diviseurs. I”, Acta Arith. 34 (1979) no. 4, p. 273-285 (loose errata)
\bibitem{DonnellyTavare87} P. Donnelly \& S. Tavaré - “The population genealogy of the infinitely-many neutral alleles model”, J. Math. Biol. 25 (1987) no. 4, p. 381-391
\bibitem{Feng10} S. Feng - The Poisson-Dirichlet distribution and related topics. Models and asymptotic behaviors, Probability and its Appl., Springer, Heidelberg, 2010
\bibitem{JKB} N. L. Johnson, S. Kotz \& N. Balakrishnan - Discrete multivariate distributions, Wiley Series in Probab. and Stat.: Applied Probab. and Stat., John Wiley \& Sons, Inc., New York, 1997
\bibitem{Kingman93} J. F. C. Kingman - Poisson processes, Oxford Studies in Probab., vol. 3, The Clarendon Press, Oxford University Press, New York, 1993
\bibitem{KotzBal00} S. Kotz, N. Balakrishnan \& N. L. Johnson - Continuous multivariate distributions. Vol. 1. Models and applications, Wiley Series in Probab. and Stat.: Applied Probab. and Stat., Wiley-Interscience, New York, 2000
\bibitem{Koukoulo19} D. Koukoulopoulos - The distribution of prime numbers, Graduate Studies in Math., vol. 203, American Mathematical Society, Providence, RI, 2019
\bibitem{LastPenrose18} G. Last \& M. Penrose - Lectures on the Poisson process, Institute of Math. Stat. Textbooks, vol. 7, Cambridge University Press, Cambridge, 2018
\bibitem{Ken22} S.-K. Leung - “Dirichlet law for factorisation of integers, polynomials and permutations”, Math. Proc. Cambridge Philos. Soc. 175 (2023) no. 3, p. 649-676
\bibitem{KenWeb24} S.-K. Leung - “Dirichlet law with arbitrary parameters for factorizations of integers” (2025), Retrieved on 24 August 2025 from https://sites.google.com/view/sunkaileung/miscellaneous
\bibitem{MV07} H. L. Montgomery \& R. C. Vaughan - Multiplicative number theory. I. Classical theory, Cambridge Studies in Advanced Math., vol. 97, Cambridge University Press, Cambridge, 2007
\bibitem{NyandwiSmati13} S. Nyandwi \& A. Smati - “Distribution laws of pairs of divisors”, Integers 13 (2013), article ID A13, 13 pages
\bibitem{Tenenbaum00} G. Tenenbaum - “A rate estimate in Billingsley’s theorem for the size distribution of large prime factors”, Q. J. Math. 51 (2000) no. 3, p. 385-403
\end{thebibliography}
\end{document}
